% Bagian: sets-functions-relations
% Bab: functions
% Subbagian: functions-relations

\documentclass[../../../include/open-logic-section]{subfiles}


\begin{document}

\olfileid[id]{sfr}{fun}{rel}

\olsection{Fungsi sebagai Relasi}

\begin{explain}
Suatu fungsi yang memetakan !!{element}s dari~$A$ ke !!{element}s dari~$B$
jelas menentukan relasi antara $A$ dan~$B$, yaitu relasi yang berlaku antara
$x$ dan $y$ jika dan hanya jika $f(x) = y$. Bahkan, jika misalnya kita ingin
mereduksi unsur-unsur dasar matematika, kita dapat \emph{mengidentifikasi}
fungsi~$f$ dengan relasi ini, yaitu dengan suatu himpunan pasangan. Hal ini
kemudian menimbulkan pertanyaan: relasi mana yang menentukan fungsi dengan
cara tersebut?
\end{explain}

\begin{defn}[Graf suatu fungsi] Misalkan $f\colon A \to B$ suatu fungsi.
\emph{Graf} dari~$f$ adalah relasi $R_f \subseteq A \times B$ yang
didefinisikan oleh
\[
R_f = \Setabs{\tuple{x,y}}{f(x) = y}.
\]
\end{defn}

\begin{explain}
Graf suatu fungsi ditentukan secara unik oleh ekstensionalitas. Selain itu,
ekstensionalitas (pada himpunan) langsung membenarkan prinsip implisit
ekstensionalitas untuk fungsi: jika $f$ dan~$g$ mempunyai domain dan kodomain
yang sama, maka keduanya identik apabila nilainya sama untuk semua argumen.

Demikian pula, jika suatu relasi bersifat ``fungsional'', relasi itu merupakan
graf suatu fungsi.
\end{explain}

\begin{prop}\ollabel{prop:graph-function}
Misalkan $R \subseteq A \times B$ sedemikian sehingga:
\begin{enumerate}
\item Jika $Rxy$ dan $Rxz$, maka $y = z$; dan
\item untuk setiap $x \in A$ terdapat suatu $y \in B$ sehingga $\tuple{x,
y} \in R$.
\end{enumerate}
Maka $R$ adalah graf fungsi $f\colon A \to B$ yang didefinisikan oleh
$f(x) = y$ jika dan hanya jika $Rxy$.
\end{prop}

\begin{proof}
Andaikan terdapat $y$ sehingga $Rxy$. Jika terdapat pula $z \neq y$ sehingga
$Rxz$, syarat pada~$R$ akan dilanggar. Jadi, jika terdapat $y$ sehingga
$Rxy$, maka $y$ tersebut unik, sehingga $f$ terdefinisi dengan baik. Jelas
bahwa $R_f = R$.
\end{proof}

\begin{explain}
Setiap fungsi $f\colon A \to B$ mempunyai graf, yaitu relasi pada $A \times B$
yang didefinisikan oleh $f(x) = y$. Sebaliknya, setiap relasi~$R \subseteq A
\times B$ dengan sifat-sifat yang diberikan dalam
\olref{prop:graph-function} merupakan graf suatu fungsi~$f \colon A \to B$.
Karena hubungan erat antara fungsi dan grafnya ini, kita dapat memandang suatu
fungsi semata-mata sebagai grafnya. Dengan kata lain, fungsi dapat
diidentifikasi dengan relasi tertentu, yaitu dengan himpunan tupel tertentu.
\oliflabeldef{sfr:rel:ref:sec}{Namun, perhatikan bahwa semangat
``identifikasi'' ini sama seperti dalam \olref[sfr][rel][ref]{sec}: ini bukan
klaim tentang metafisika fungsi, melainkan pengamatan bahwa kita mudah
\emph{memperlakukan} fungsi sebagai himpunan tertentu. Salah satu alasan
perlakuan ini begitu memudahkan ialah bahwa k}{K}ita sekarang dapat meninjau
operasi pada fungsi yang serupa dengan operasi yang kita lakukan pada relasi
(lihat \olref[sfr][rel][ops]{sec}). Secara khusus:
\end{explain}

\begin{defn}\ollabel{defn:funimage}
Misalkan $f \colon A \to B$ suatu fungsi dengan $C\subseteq A$.

\emph{Pembatasan} dari~$f$ pada~$C$ adalah fungsi
$\funrestrictionto{f}{C}\colon C \to B$ yang didefinisikan oleh
$(\funrestrictionto{f}{C})(x) = f(x)$ untuk semua $x \in C$. Dengan kata lain,
$\funrestrictionto{f}{C} = \Setabs{\tuple{x, y} \in R_f}{x \in C}$.

\emph{Penerapan} dari~$f$ pada~$C$ adalah $\funimage{f}{C} =
\Setabs{f(x)}{x \in C}$. Kita juga menyebutnya \emph{peta} dari~$C$ di bawah
fungsi~$f$.
\end{defn}

\begin{explain}
Dari definisi-definisi ini diperoleh $\ran{f} =
\funimage{f}{\dom{f}}$, untuk setiap fungsi~$f$.
\oliflabeldef{sfr:rel:ops:sec}{Gagasan-gagasan ini tepat seperti yang kita
harapkan berdasarkan definisi dalam \olref[sfr][rel][ops]{sec} dan
identifikasi fungsi dengan relasi. Namun, dua operasi lain---invers dan hasil
kali relatif---memerlukan sedikit penjelasan tambahan. Penjelasan itu akan
kita berikan dalam \olref[inv]{sec} dan \olref[cmp]{sec}.}{}
\end{explain}

\end{document}
